\reading{CR-5}{Introduction to Credit Risk Modeling and Assessment}{STRIKE FIRST}{8}

\begin{crkeybox}[Learning objectives for this reading]
\begin{enumerate}[label=\textbf{\alph*.},leftmargin=2em]
\item Explain the capital adequacy, asset quality, management, earnings, and liquidity (CAMEL) system used for evaluating the financial condition of a bank.
\item Describe quantitative measurements and factors of credit risk, including probability of default, loss given default, exposure at default, expected loss, and time horizon.
\item Estimate risk-weighted assets and capital adequacy ratio of a financial institution.
\item Describe the judgmental approaches, empirical models, and financial models to predict default.
\item Apply the Merton model to calculate default probability and the distance to default and describe the limitations of using the Merton model.
\item Compare and contrast different approaches to credit risk modeling, such as those related to the Merton model, Credit Risk Plus (CreditRisk+), CreditMetrics, and the Moody's-KMV model.
\item Apply risk-adjusted return on capital (RAROC) to measure the performance of a loan.
\end{enumerate}
\end{crkeybox}

% =====================================================================
\lo{CR-5 a}{Explain the capital adequacy, asset quality, management, earnings, and liquidity (CAMEL) system used for evaluating the financial condition of a bank.}

Analytical tools for credit risk management are widely used. One of the most
common is \term{CAMEL} (capital, assets, management, earnings, and liquidity),
which analyses a bank's overall health using five factors.

\begin{enumerate}
\item \term{Capital adequacy.} Regulation establishes a minimum capital reserve
that all banks must maintain. This analysis should consider historical and
current positioning. \emph{How much reserve does the bank have to cover losses?}
\item \term{Asset quality.} Reviewers examine whether the bank's loans are
performing or showing some signs of delinquency. This analysis considers interest
rate risk and liquidity risk in relation to the bank's assets. \emph{Are the
bank's loans performing well, or are there defaults?}
\item \term{Management.} This portion of the analysis focuses on the quality of
the bank's risk management structure. Business strategy, financial performance,
and internal controls and policies are all considered. \emph{Does the bank have
strong risk management practices?}
\item \term{Earnings.} Core earnings should be reviewed. Examiners are
specifically looking for the stability of earnings, net interest margin, return on
assets, and future earnings potential, with an eye for the anticipated impact of
stressed economic events. \emph{Is the bank profitable and prepared for economic
downturns?}
\item \term{Liquidity.} This category considers both the impact of interest rate
risk on the bank's investments and the influence of liquidity risk. The goal is to
understand if the bank has the potential for short-term liquidity events that
could jeopardize the bank's viability. \emph{Does the bank have enough cash to
meet short-term obligations?}
\end{enumerate}

% =====================================================================
\lo{CR-5 b}{Describe quantitative measurements and factors of credit risk, including probability of default, loss given default, exposure at default, expected loss, and time horizon.}

\begin{enumerate}
\item \term{Probability of default (PD).} This captures the chance of a borrower
defaulting on their loan, typically defined as a late payment exceeding a certain
threshold (e.g., 90 days). It is expressed as a percentage.
\item \term{Exposure at default (EAD).} This represents the outstanding amount
owed on the loan at the time of default. It is influenced by loan characteristics
and can be a fixed amount or fluctuate (like credit lines). EAD is expressed in
dollars.
\item \term{Loss given default (LGD).} This reflects the percentage of the loan
amount that is expected to be unrecoverable after default. Recovery rate (RR) is
the inverse, calculated as 100\% minus LGD. Borrower-specific factors like
industry and financial health can influence LGD estimates.
\end{enumerate}

\begin{center}
\begin{tikzpicture}[font=\sffamily\footnotesize]
\node[draw=FRMGreen, fill=PaleGreen, line width=0.8pt, rounded corners=1pt,
      text width=2.9cm, align=center, minimum height=1.25cm] (el) at (0,0)
      {\textbf{Expected loss}\\\scriptsize (EL)\\\scriptsize \$};
\node[font=\large] (eq) at (1.95,0) {$=$};
\node[draw=FRMNavy, fill=PaleNavy, line width=0.8pt, rounded corners=1pt,
      text width=2.9cm, align=center, minimum height=1.25cm] (pd) at (3.9,0)
      {\textbf{Probability of default}\\\scriptsize (PD)\\\scriptsize \%};
\node[font=\large] (x1) at (5.85,0) {$\times$};
\node[draw=FRMGold, fill=PaleGold, line width=0.8pt, rounded corners=1pt,
      text width=2.9cm, align=center, minimum height=1.25cm] (lgd) at (7.8,0)
      {\textbf{Loss given default}\\\scriptsize (LGD)\\\scriptsize \%};
\node[font=\large] (x2) at (9.75,0) {$\times$};
\node[draw=FRMGold, fill=PaleGold, line width=0.8pt, rounded corners=1pt,
      text width=2.9cm, align=center, minimum height=1.25cm] (ead) at (11.7,0)
      {\textbf{Exposure at default}\\\scriptsize (EAD)\\\scriptsize \$};
% brackets
\draw[FRMNavy, line width=0.8pt] (2.45,1.0) -- (2.45,1.3) -- (5.35,1.3) -- (5.35,1.0);
\node[above, font=\scriptsize\bfseries, color=FRMNavy] at (3.9,1.3)
  {Borrower risk quantification};
\draw[FRMGold, line width=0.8pt] (6.35,1.0) -- (6.35,1.3) -- (13.15,1.3) -- (13.15,1.0);
\node[above, font=\scriptsize\bfseries, text=FRMGold] at (9.75,1.3)
  {Facility risk quantification};
\node[below, font=\scriptsize, color=FRMGrey, align=left] at (6.5,-0.85)
  {Inputs: risk rating \ \textbullet\ \ type of counterparty \ \textbullet\ \ guarantees};
\end{tikzpicture}
\end{center}
\figcap{Expected loss calculation. The split above the boxes is the fact worth
holding: PD is a property of the \emph{borrower}, while LGD and EAD are
properties of the \emph{facility}. Two obligations to the same borrower share a
PD but can carry very different LGD and EAD.}

\gapbadge

{\footnotesize\color{FRMGrey}The source notes cover PD, EAD and LGD but do not
treat the time horizon, which the objective names explicitly. The following is
written from the GARP curriculum.}

\begin{crgapbox}[What the objective asks for: time horizon]
The main outcome of credit risk modelling and management systems is an estimate
of the expected loss for a loan \emph{over a specific period}, which is usually
set to be one year, in accordance with the standard time basis used for financial
reporting.

\smallskip
Horizon also enters as a risk factor in its own right: the duration of the loan
matters, because short-term allocations are generally less risky than long-term
liabilities, due to the increasing uncertainty over longer time horizons.
\end{crgapbox}

% =====================================================================
\lo{CR-5 c}{Estimate risk-weighted assets and capital adequacy ratio of a financial institution.}

Banks should always have enough capital to match the risk of their assets. This
amount of capital can be measured using the \term{capital adequacy ratio (CAR)}.
A higher CAR indicates that a bank has a larger capital buffer and is therefore
better equipped to absorb losses.

\begin{crfmlbox}[Capital adequacy ratio]
\[ \mathrm{CAR} \;=\; \frac{\text{capital}}{\text{risk-weighted assets}} \;\geq\; \alpha \]
\vk{capital $=$ Tier 1 and Tier 2 capital; $\alpha$ $=$ the minimum requirement
imposed by the regulator, 8\% under Basel II and 10.5\% under Basel III.}
\end{crfmlbox}

\term{Risk-weighted assets (RWA)} take into account the riskiness of a bank's
assets. Assets are assigned a risk weight depending on the type of borrower and
riskiness. There are two methods for deriving a bank's RWA.

\subsection*{1) Standardized approach}

A supervisory authority can prescribe predetermined weights for various asset
categories. This method is easy to deploy, but it is based on data that is
external to a given bank, so it is easy but less precise.

\begin{center}
\footnotesize
\begin{tabular}{L{2.2cm} L{12.3cm}}
\toprule
\thead{Risk weight} & \thead{Asset category} \\
\midrule
0\% & Cash, gold, claims on Organization of Economic Co-operation and
Development (OECD) countries such as U.S. Treasury bonds, and insured residential
mortgages \\
20\% & Claims on OECD banks and government agencies like U.S. agency securities
such as Fannie Mae and Freddie Mac, or municipal bonds \\
50\% & Uninsured residential mortgages \\
100\% & Loans to corporations, consumer loans, corporate bonds, claims on
non-OECD banks \\
\bottomrule
\end{tabular}
\end{center}
\figcap{Standardized-approach risk weights. Note that the insurance status of a
residential mortgage moves it between the 0\% and the 50\% bucket, and that OECD
membership is what separates the 20\% bucket from the 100\% one.}

\subsection*{2) Internal ratings-based (IRB) approach}

The IRB approach is customized to each bank and requires sophisticated modelling
using historical data. Internal credit scores can be combined with historical
data to estimate the PD and LGD.

The RWA is then calculated using an \term{asymptotic single risk factor (ASRF)}
model, which assumes the credit portfolio is diversified such that no
idiosyncratic risk is left, that is, only systematic risk remains.

The maturity adjustment has two components:

\begin{enumerate}[label=\alph*)]
\item First, capital requirements should \emph{increase} with maturity, because
the risk of default increases as maturity increases.
\item Second, borrowers with \emph{low} probabilities of default should have
\emph{higher} adjustments than high PD borrowers. A borrower with a high PD
already has a higher instance of default factored into their estimates.
\end{enumerate}

Basel III also strengthened bank capital requirements by:

\begin{itemize}
\item increasing the minimum capital ratio from 8\% to 10.5\%;
\item adding rules on managing debt (leverage) and access to cash (liquidity);
\item focusing on the risk of defaults by other banks (counterparty risk) in
complex financial products; and
\item requiring stricter risk management practices like stress testing and model
validation.
\end{itemize}

% =====================================================================
\lo{CR-5 d}{Describe the judgmental approaches, empirical models, and financial models to predict default.}

There are three primary approaches for assessing the likelihood of default in
credit risk models.

\subsection*{1) Judgmental approach --- usable for both corporate and consumer loans}

Relies on the experience and expertise of credit risk officers to assess the
creditworthiness of a borrower. Factors known as the \term{5 Cs} --- the
borrower's character, capacity, capital, collateral, and economic conditions ---
are considered.

\begin{itemize}
\item \term{Advantages.} Useful when historical data is scarce (e.g., project
finance where historical data may not be available); leverages the experience of
credit analysts and benefits from human insight; flexible in assessing unique
loan types.
\item \term{Disadvantages.} Output quality is hard to validate, since human
judgements are involved; may not update swiftly to changing conditions; impact on
the entire credit portfolio is difficult to assess; lacks transparency and
consistency.
\end{itemize}

\subsection*{2) Empirical (data-driven) models --- usable for both corporate and consumer loans}

Use historical data (e.g., loans accepted, rejected, in good standing, and in
default) to search for nontrivial patterns in the relationship between the
likelihood of default and input variables, that is, risk factors. This
data-driven approach can harness the power of machine learning (ML) to search for
highly complex, nonlinear relationships.

\begin{itemize}
\item \term{Advantages.} Transparency and consistency in the credit risk process;
predictive power and model structure can be validated empirically; ability to
explore hypothesized relationships analytically; real-time market data updates to
adjust to changing conditions; sector-specific market adjustments are possible.
\item \term{Disadvantages.} Past data may not predict future outcomes, especially
in volatile conditions; some inputs like financial statements are not updated in
real time, potentially delaying risk information.
\end{itemize}

\subsection*{3) Theory-based financial models (market models) --- corporate loans only}

Rely on economic and financial theory and financial market data. There are two
primary types.

\begin{enumerate}[label=\alph*)]
\item \term{Structural models.} This method assumes that default is an
\emph{endogenous} process, which is influenced by the firm's internal structural
characteristics (e.g., the Merton model).
\item \term{Reduced-form models.} This method assumes that default is driven by a
random \emph{exogenous} event (e.g., a Poisson jump process). The primary input
sources for the credit risk structure are market data on credit derivatives and
bonds.
\end{enumerate}

\begin{center}
\footnotesize
\begin{tabular}{L{2.1cm} L{4.1cm} L{4.1cm} L{4.1cm}}
\toprule
 & \thead{Judgmental approach} & \thead{Empirical models} & \thead{Financial models} \\
\midrule
Data source & Qualitative data from experts (credit risk officers) &
  Historical data from internal databases, CRAs, and external sources &
  Market data on credit derivatives and bonds \\
Methodology & Relies on expert judgment and qualitative analysis &
  Analyzes historical data to identify patterns and relationships &
  Based on normative economic and financial theory \\
Inputs & 5C analysis: character, capacity, capital, collateral, conditions &
  Borrower characteristics, loan status, external risk factors &
  Market data on credit derivatives, bonds, and financial indicators \\
Applicability & Can be used for both consumer and corporate loans &
  Applicable to various loan types, consumer and corporate &
  Primarily used for corporate borrowers \\
Strengths & Useful for unique loans; provides qualitative insights &
  Offers transparency, empirical validation, and real-time updates &
  Based on economic theory; provides market-based assessments \\
Weaknesses & Relies on expert judgment; lacks transparency and consistency &
  May suffer during uncertain periods; requires accurate historical data &
  Limited to corporate borrowers; may not capture market changes quickly \\
\bottomrule
\end{tabular}
\end{center}
\figcap{The three approaches side by side. Applicability is the row that
discriminates fastest: only the financial models are restricted to corporate
borrowers.}

% =====================================================================
\lo{CR-5 e}{Apply the Merton model to calculate default probability and the distance to default and describe the limitations of using the Merton model.}

The Merton model is a \term{structural} approach to predicting default. It
assumes default occurs when the proprietary structure of the defaulting company is
no longer worthwhile. It is based on Black-Scholes-Merton option pricing theory
and evaluates various components of firm value.

\begin{crkeybox}[Notation used in this section]
\begin{tabular}{L{1.9cm} L{12.6cm}}
$V_0$ & current value of the firm's assets \\
$V_T$ & value of the firm's assets at time $T$ \\
$D$ & total debt to be repaid at time $T$ \\
$E_0$ & current value of the firm's equity \\
$E_T$ & value of the firm's equity at time $T$ (i.e., $V_T - D$) \\
$\sigma_V$ & asset volatility \\
$\sigma_E$ & equity volatility \\
\end{tabular}
\end{crkeybox}

\begin{crtrapbox}[Three symbols for the same object]
The source notes write the face value of debt as \term{F} in the payoff formulas,
as \term{D} in the Black-Scholes-Merton equity formula, and as \term{L} in the
real-world probability of default formula. All three denote the same thing: the
amount owed at $T$, which is the strike price of the option. They are carried
here as written rather than silently harmonised, because the exam material draws
on all three sources.
\end{crtrapbox}

\subsection*{Equity as a call, debt as a short put}

\begin{crfmlbox}[Terminal payoffs]
\[ S_T \;=\; \max(V_T - F,\,0) \qquad\qquad D_T \;=\; F - \max(F - V_T,\,0) \]
\vk{$S_T$ = value of equity at $T$; $D_T$ = value of debt at $T$; $F$ = principal
amount due on the zero-coupon bond. The equity payoff is that of a \emph{long
call option} on the firm's assets struck at $F$. The debt payoff is the face
amount less a \emph{short put} struck at $F$.}
\end{crfmlbox}

\begin{center}
\begin{tikzpicture}
\begin{axis}[
  width=13.0cm, height=5.8cm,
  axis lines=left,
  xlabel={Firm asset value at $T$, \ $V_T$ (\$m)},
  ylabel={Value at $T$ (\$m)},
  xlabel style={font=\sffamily\scriptsize}, ylabel style={font=\sffamily\scriptsize},
  tick label style={font=\sffamily\scriptsize},
  xmin=0, xmax=100, ymin=0, ymax=68,
  xtick={0,25,50,75,100}, ytick={0,25,50},
  legend style={font=\scriptsize\sffamily, draw=FRMRule, fill=white,
    at={(0.5,-0.32)}, anchor=north, legend columns=-1, column sep=10pt}]
\addplot[draw=FRMNavy, line width=1.2pt, sharp plot]
  coordinates {(0,0) (50,0) (100,50)};
\addlegendentry{equity $S_T=\max(V_T-F,0)$}
\addplot[draw=FRMGold, line width=1.2pt, sharp plot]
  coordinates {(0,0) (50,50) (100,50)};
\addlegendentry{debt $D_T=F-\max(F-V_T,0)$}
\draw[FRMRed, dashed, line width=0.8pt] (axis cs:50,0) -- (axis cs:50,58);
\node[font=\sffamily\scriptsize\bfseries, text=FRMRed, fill=white, inner sep=1.5pt]
  at (axis cs:50,62) {$F=50$};
\addplot[only marks, mark=*, mark size=1.7pt, draw=FRMGreen, fill=FRMGreen]
  coordinates {(60,10) (40,0)};
\node[font=\sffamily\scriptsize, text=FRMGreen, fill=white, inner sep=1.5pt]
  at (axis cs:74,13) {$V_T=60 \Rightarrow S_T=10$};
\node[font=\sffamily\scriptsize, text=FRMGreen, fill=white, inner sep=1.5pt]
  at (axis cs:24,7) {$V_T=40 \Rightarrow S_T=0$};
\end{axis}
\end{tikzpicture}
\end{center}
\figcap{The firm's asset value at $T$ is split between the two claimants, and the
split is a kinked function with its kink at the face value of the debt. Below
$F$ the debtholders take everything and the equity is worth nothing; above $F$
the debtholders are capped at $F$ and every further dollar goes to equity. The
two worked points are the examples below.}

\begin{crexbox}[Equity value at $T$]
Calculate the value of the firm's equity at $T$, $S_T$, given that the principal
amount due on the zero-coupon bond is \$50 million and the total value of the firm
at $T$, $V_T$, is \$60 million. In addition, what is the value of equity if $V_T$
is \$40 million?
\tcblower
\[ \text{Equity value} = \max(60 - 50,\,0) = 10 = \$10\ \text{million} \]
\[ \text{Equity value} = \max(40 - 50,\,0) = 0 = \$0\ \text{million} \]
\end{crexbox}

\begin{crexbox}[Debt value at $T$]
Calculate the value of the firm's debt at $T$, $D_T$, given that the principal
amount due on the zero-coupon bond is \$50 million and the total value of the firm
at $T$, $V_T$, is \$40 million. In addition, what is the value of debt if $V_T$ is
\$60 million?
\tcblower
\[ \text{Debt value} = 50 - \max(50 - 40,\,0) = 50 - 10 = 40 = \$40\ \text{million} \]
\[ \text{Debt value} = 50 - \max(50 - 60,\,0) = 50 - 0 = 50 = \$50\ \text{million} \]
\end{crexbox}

\subsection*{Value of equity using the BSM model}

On the assumption of no dividends, the Black-Scholes-Merton equity value today is
stated as follows.

\begin{crfmlbox}[Equity as a call on firm assets]
\[ E_0 \;=\; V_0 N(d_1) \;-\; D e^{-rT} N(d_2) \]
\[ d_1 = \frac{\ln(V_0/D) + \bigl(r + \sigma_V^{2}/2\bigr)T}{\sigma_V\sqrt{T}}
   \qquad\qquad
   d_2 = d_1 - \sigma_V\sqrt{T} \]
\vk{$N(\cdot)$ = cumulative normal distribution function; $r$ = risk-free rate;
$T$ = time to repayment in years. Note the sign inside $d_1$: the drift term
carries $+\sigma_V^{2}/2$, while $d_2$ written out in full carries
$-\sigma_V^{2}/2$.}
\end{crfmlbox}

\subsection*{Distance to default}

Distance to default measures the distance between a firm's current financial
position and the point of default. DtD is the number of standard deviations the
company's assets are away from the face value of the debt.

\begin{crfmlbox}[Distance to default]
\[ \mathrm{DtD} \;=\; d_2 \;=\;
   \frac{\ln(V_0) - \ln(D) + \bigl(r - \sigma_V^{2}/2\bigr)T}{\sigma_V\sqrt{T}} \]
\end{crfmlbox}

\begin{crkeybox}[Direction]
The higher (lower) the distance to default, the lower (higher) the probability of
default.
\end{crkeybox}

\subsection*{Probability of default}

\begin{itemize}
\item The \term{risk-neutral probability of default}, $N(-d_2)$, is effectively
the probability that shareholders will not exercise their option to repay the
loan. This analysis assumes that the firm's assets grow at the risk-free rate.
\item For the \term{real-world probability of default}, the expected return on the
firm's assets $(\mu)$ is substituted for the risk-free rate $(r)$.
\end{itemize}

\begin{crfmlbox}[The two probabilities of default]
\[ \text{Risk-neutral PD} \;=\; N(-d_2), \qquad N(-d_2) = 1 - N(d_2) \]
\[ \text{Real-world PD} \;=\;
   N\left[\,-\,\frac{\ln\!\left(\dfrac{A}{L}\right) + \left(\mu - \dfrac{\sigma_A^{2}}{2}\right)T}
   {\sigma_A\sqrt{T}}\,\right] \]
\vk{$A$ = current market value of the firm's assets; $L$ = face value of the debt
(the default point); $\mu$ = expected return on the firm's assets; $\sigma_A$ =
asset volatility. The only change from the risk-neutral version is $r$ becoming
$\mu$.}
\end{crfmlbox}

\gapbadge

{\footnotesize\color{FRMGrey}The objective asks for the limitations of the Merton
model. The source notes carry the objective heading but no body text for this
half. The following is written from the GARP curriculum.}

\begin{crgapbox}[What the objective asks for: limitations of the Merton model]
\begin{enumerate}
\item \term{The capital structure is assumed to be trivially simple.} The
derivation assumes a firm with a single liability of face value $L$ maturing at a
single date $T$. In practice $L$ is usually taken as the short-term debt reported
on the balance sheet, which is a choice, not an observation.
\item \term{The two key inputs are unobservable.} The current value of the
firm's assets $A$ and the asset volatility $\sigma_A$ cannot be read off a
market. They must be backed out by solving the equity-value equation and the
equity-volatility equation simultaneously, which requires an iterative
Newton-type algorithm.
\item \term{The mapping from distance to default to a probability relies on
normality.} The model converts DtD into a PD through the normal distribution. The
Moody's-KMV EDF model exists precisely because that mapping is unrealistic: it
replaces the normal distribution with an empirical distribution of default
frequencies derived from historical data.
\item \term{The real-world PD needs $\mu$, which must itself be estimated.} The
standard approach is to solve the system first and then define $\mu$ as the rate
of change in $A$ over past time periods, so the real-world probability inherits
the estimation error of that step.
\end{enumerate}
\end{crgapbox}

% =====================================================================
\lo{CR-5 f}{Compare and contrast different approaches to credit risk modeling, such as those related to the Merton model, Credit Risk Plus (CreditRisk+), CreditMetrics, and the Moody's-KMV model.}

\begin{enumerate}
\item \term{Merton model.} Uses an option pricing model and the borrower's
short-term debt to estimate PD. Assumes a standardized normal distribution for
default frequencies. The default point $L$ considers only short-term debt.
\item \term{Moody's-KMV Expected Default Frequency (EDF) model.} Differs from
Merton in two ways: it uses historical data to develop a distribution of default
frequencies, which is more realistic; and the default point $L$ includes
short-term debt \emph{and half of long-term debt}, which better reflects loan
obligations.
\item \term{CreditMetrics model.} Primarily used for estimating the market value
of nonmarketable loans or loan portfolios, employing a peer-set analysis
approach. This involves comparing a borrower's default risk to that of others,
using a rating transition matrix that reflects changes in default rates over time.
Model inputs include the borrower's credit rating, a transition matrix (detailing
credit rating changes and default rates), historical recovery rates, and the bond
market's yield margins.
\end{enumerate}

\gapbadge

{\footnotesize\color{FRMGrey}The objective names four models. The source notes
cover three and omit Credit Risk Plus entirely. The following is written from the
GARP curriculum. CreditRisk+ is treated in full at CR-10.}

\begin{crgapbox}[What the objective asks for: Credit Risk Plus (CreditRisk+)]
In 1997 Credit Suisse Financial Products proposed a methodology for calculating
VaR that it termed Credit Risk Plus. It involves analytic approximations that are
well established in the \emph{insurance} industry.

\smallskip
Suppose an institution has $n$ loans of a certain type and the probability of
default for each loan during a year is $q$, so the expected number of defaults is
$qn$. Assuming default events are independent, the probability of $m$ defaults is
binomial. If $q$ is small and $n$ large, this is approximated by the Poisson
distribution:
\[ \text{Prob}(m \text{ defaults}) \;=\; \frac{e^{-qn}(qn)^{m}}{m!} \]
This holds approximately even when the probability of default is not the same for
each loan, provided all the default probabilities are small and $q$ is the average
probability of default over the next year for the loans in the portfolio.

\smallskip
In practice the default rate $q$ is itself uncertain and varies greatly from year
to year. A convenient assumption is that $qn$ follows a gamma distribution, in
which case the Poisson distribution becomes a \emph{negative binomial}
distribution. As the standard deviation of $qn$ tends to zero the negative
binomial converges back to the Poisson; as it increases, the probability of an
extreme outcome involving a large number of defaults rises.
\end{crgapbox}

\begin{center}
\footnotesize
\begin{tabular}{L{2.4cm} L{3.7cm} L{3.9cm} L{4.4cm}}
\toprule
 & \thead{Merton} & \thead{Moody's-KMV EDF} & \thead{CreditMetrics} \\
\midrule
Default point $L$ & Short-term debt only &
  Short-term debt $+$ half of long-term debt &
  Not a default-point model; uses a rating transition matrix \\
Distribution used & Standardized normal &
  Empirical distribution of historical default frequencies &
  Transition probabilities from the migration matrix \\
Primary use & PD for a single firm from option pricing &
  PD for a single firm, calibrated to history &
  Market value of nonmarketable loans and loan portfolios \\
\bottomrule
\end{tabular}
\end{center}
\figcap{Merton versus Moody's-KMV is the sibling pair that matters most here.
They share the option-pricing skeleton and differ on exactly two things: what
goes into the default point, and whether the distance to default is mapped to a
probability through the normal distribution or through history.}

% =====================================================================
\lo{CR-5 g}{Apply risk-adjusted return on capital (RAROC) to measure the performance of a loan.}

RAROC is a method used to assess the profitability of a business considering the
risk involved. It considers the income generated by a loan relative to the capital
required to cover the default risk: a metric used to assess the performance of a
loan by comparing financial outcomes to the economic capital required for those
outcomes.

\begin{itemize}
\item \term{Loan revenue factors.} The spread $(s)$ between the loan rate and the
bank's cost of capital; the fees $(f)$ attributable to the loan; the estimated
loan losses $(l)$; the associated operating costs $(c)$; and relevant taxation
$(x)$.
\item \term{Capital at risk.} The loan's value $(L)$, its duration $(D)$, the
current interest rate $(i)$, and the estimated change in interest rates
$(\Delta i)$.
\end{itemize}

\begin{crfmlbox}[RAROC]
\[ \mathrm{RAROC} \;=\; \frac{\text{loan revenues}}{\text{capital at risk}} \]
\[ \text{loan revenues} \;=\; \text{loan value} \times (s + f - l - c)(1-x) \]
\[ \Delta L \;\approx\; -\,L \times D \times \frac{\Delta i}{1+i} \]
\vk{$s$ = spread, $f$ = fees, $l$ = expected loan losses, $c$ = operating costs,
$x$ = tax rate, all as decimals; $L$ = loan value; $D$ = duration; $i$ = current
interest rate; $\Delta i$ = expected change in rates. Note that the tax factor
$(1-x)$ multiplies the \emph{whole} bracket, and that $l$ and $c$ are subtracted
while $s$ and $f$ are added.}
\end{crfmlbox}

\begin{crexbox}[XYZ Bank: RAROC on a loan]
Assume that XYZ Bank has a loan with a value of \$1,000,000. There is an
associated commission of 0.1\%, the spread between the loan's interest rate and
the bank's cost of capital is 0.65\%, the estimate of operating costs is 0.2\%,
and their estimated tax rate is 11.1\%. Additionally, this loan has a duration of
3.75 and an interest rate of 10.5\%. Interest rates are expected to increase by
1.5\%.
\tcblower
\textbf{Loan revenues.}
\[ \$1{,}000{,}000\,(0.0065 + 0.001 - 0 - 0.002)\,(1 - 0.111) = \$4{,}889.50 \]
The bracket is $0.0055$; the after-tax factor is $0.889$; the product with
\$1,000,000 is \$4,889.50. Note that estimated loan losses $l$ are zero here
because none were given.

\textbf{Capital at risk.}
\[ \Delta L = -\,\$1{,}000{,}000\,(3.75)\left(\frac{0.015}{1.105}\right) = -\$50{,}904.98 \]
The rate change is divided by $1+i = 1.105$, not by $i$.

\textbf{RAROC.}
\[ \mathrm{RAROC} = \frac{\$4{,}889.50}{\$50{,}904.98} = 9.61\% \]
\end{crexbox}

\begin{crkeybox}[Decision rule]
The loan is profitable if its RAROC exceeds the \term{bank's cost of capital}.
\end{crkeybox}

\subsection*{An alternative denominator}

There is an alternative method for estimating the capital at risk that uses
historical values, and not market-based metrics like expected changes in interest
rates. This method uses the unexpected loan loss in the RAROC denominator.

\begin{crfmlbox}[Unexpected loan loss as capital at risk]
\[ \text{unexpected loan loss} \;=\; \alpha \times \mathrm{LGD} \times \mathrm{EAD} \]
\vk{$\alpha$ = an adjustment for unexpected default rates, derived from the
distribution of historical default rates; LGD = loss given default; EAD =
exposure at default.}
\end{crfmlbox}

If the distribution of default rates is assumed to be normal, then at the 99.5\%
confidence level an analyst should set $\alpha = 2.6\sigma$. However, most loan
distributions have some degree of skewness. For this reason it is more common to
use a higher value (e.g., 5 or 6) based on the characteristics of the loan.
